% Author: OUYANG Di. Insert into the group's CVPR document.
% Requires no packages beyond the CVPR template's standard math support.
% This is a methods fragment, not a standalone document or an empirical result.

\subsection{Data loading and common preprocessing}
\label{sec:ouyang-data}
We load images from the frozen manifest using its split, relative path, and
binary label, where real and AI-generated images are labelled 0 and 1,
respectively. The image identifier and generator metadata are retained for
prediction alignment and audit. We keep the training, validation, seen-test,
ADM-test, and GLIDE-test partitions separate; the development fixture is not
added to the formal training or evaluation sets. These partitions inherit the
mirror dataset's frozen re-split of the upstream GenImage training partition,
rather than the original GenImage training/validation directory split.

All conditions use the same preprocessing. We correct EXIF orientation,
convert to RGB, resize the shorter side to 256 pixels with bicubic
interpolation while preserving aspect ratio, and take a $224\times224$ centre
crop. Pixel intensities are scaled to $[0,1]$ and normalised using ImageNet
channel means $(0.485,0.456,0.406)$ and standard deviations
$(0.229,0.224,0.225)$. The clean input receives no additional corruption;
compression or resampling already present in the source image is retained.

\subsection{Paired views and reproducible augmentation}
\label{sec:ouyang-views}
For conditions B and C, each sample supplies the common clean crop and a
transformed view of the same crop. We draw uniformly from three modes:
JPEG recompression, downsampling followed by resizing back to the crop size,
or downsampling followed by JPEG recompression. JPEG quality is sampled
uniformly from the integers 50 through 95; the downsampling factor is sampled
uniformly between 0.5 and 0.9. Both labels follow the same transformation
distribution. Corruptions operate on RGB pixels before tensor conversion and
normalisation. No additional random crop, flip, or colour perturbation is
introduced by this module. Resizing down and back uses bicubic interpolation;
downsampled dimensions are floored and bounded below by one pixel. JPEG
encoding uses Pillow with 4:2:0 subsampling, without progressive encoding or
optimisation. Exact codec and dependency versions are retained in run records.

Transformation parameters are determined by a stable seed derived from the
shared run seed, epoch, and image identifier. The derivation excludes the
condition, label, worker identifier, and access order. Consequently, B and C
receive identical views when their run seed, epoch, manifest, and data
configuration agree. An epoch-dependent sampler also reproduces the batch
order when loader settings agree. The training procedure must separately
match optimisation and update budgets; matching views alone does not enforce
those controls. Run records should retain the code revision, dependency
versions, manifest and configuration hashes, seed, epoch, and loader settings.

\subsection{Fixed robustness tests and bias control}
\label{sec:ouyang-robustness}
For each of the seen, ADM, and GLIDE test sets, the fixed suites comprise
clean preprocessing, JPEG quality 70, JPEG quality 50, $0.5\times$
downsampling, and Gaussian blur. JPEG recompression is applied to the common
$224\times224$ crop and decoded before normalisation. Downsampling changes
the crop from $224\times224$ to $112\times112$ and back. Blur uses Pillow's
GaussianBlur implementation with radius 1.0. These suites preserve image IDs
so that clean and corrupted predictions can be paired. The classification
threshold is fixed at 0.5, and checkpoint selection uses validation data only.

The source dataset has format and resolution confounds: real images are JPEG
with varying dimensions, whereas generated images are PNG with
generator-specific dimensions. We therefore provide a separate input-control
suite that resizes every RGB image to $256\times256$, recompresses it as JPEG
at quality 95, decodes it, and then applies the common preprocessing. The
control runs in memory without overwriting the source files or manifest.
It standardises the newly applied encoding and dimensions, but cannot remove
traces of previous JPEG compression or native-resolution history. Its
results should therefore be reported as a sensitivity check, rather than
proof that shortcut learning has been eliminated.

\subsection{Data reception and leakage checks}
\label{sec:ouyang-qa}
Reception QA checks manifest consistency, unique identifiers and paths,
label--generator agreement, held-out generator isolation, and SHA256/source
MD5 overlap across the formal partitions and development fixture. Full-image
checks compare actual file hashes, decoding, dimensions, format, and file
size against the manifest. Perceptual hashing retrieves cross-partition
near-duplicate candidates for manual review; perceptual similarity alone is
not treated as a duplicate label and does not trigger automatic deletion.
Any accepted partition correction requires a new manifest version and
corresponding updates to the training records. Missing files, unperformed
perceptual checks, and unresolved candidates must remain explicit in the QA
record. This methods description does not assert that the current reception
audit or manual review has passed; empirical audit outcomes belong in the
final data-quality record.
